\setauthor{AC}
\section{Motivation und Problemstellung der Abo-Verwaltung}
Die Nutzung digitaler Streaming-Dienste ist in den vergangenen Jahren zu einem festen Bestandteil des Medienalltags geworden. Viele Nutzende verwenden heute nicht nur einen einzelnen Anbieter, sondern kombinieren mehrere Plattformen parallel, um auf unterschiedliche Filmangebote zugreifen zu können. Dadurch steigt jedoch auch die Komplexität der persönlichen Medienverwaltung. Laufende Kosten verteilen sich auf mehrere Dienste, Kündigungs- und Verlängerungszeitpunkte geraten leicht aus dem Blick und es ist oft unklar, welche Inhalte bereits durch bestehende Abonnements abgedeckt sind.

Aus dieser Situation ergibt sich ein konkretes praktisches Problem: Zwar existieren Plattformen zur Anzeige von Filminformationen und Anbieter-Verfügbarkeiten, ebenso wie allgemeine Werkzeuge zur Verwaltung von Zahlungen oder Abonnements, eine zentrale Anwendung zur Verbindung dieser beiden Bereiche ist jedoch deutlich seltener. Genau an dieser Stelle setzt das Projekt \glqq Mediahub\grqq{} an. Ziel ist es, Filmrecherche, Anbieterübersicht und persönliche Abo-Verwaltung in einer gemeinsamen Anwendung zusammenzuführen und dadurch einen realen Mehrwert für Nutzende zu schaffen.

Im Rahmen der Projektarbeit wurden dafür sowohl technische als auch konzeptionelle Fragestellungen betrachtet. Dazu zählen unter anderem die Ausrichtung auf eine geeignete Zielgruppe, die Frage nach relevanten Streaming-Anbietern, die Orientierung an bestehenden Benutzeroberflächen aus verwandten Anwendungsbereichen sowie die sichere Verwaltung geschützter Benutzerfunktionen. Die Einleitung bildet damit den Ausgangspunkt für die fachliche und technische Betrachtung der gesamten Diplomarbeit.

\section{Zielsetzung der Diplomarbeit}
Das zentrale Ziel der Diplomarbeit besteht in der Konzeption und Umsetzung einer Webanwendung zur zentralen Verwaltung von Streaming-Abonnements mit Fokus auf Filme. Die Anwendung soll Nutzende dabei unterstützen, ihre aktiven Dienste übersichtlich zu verwalten, Abrechnungsinformationen nachzuvollziehen und gleichzeitig schnell zu erkennen, auf welchen Plattformen ein gesuchter Film verfügbar ist. Dadurch soll nicht nur der organisatorische Aufwand reduziert, sondern auch eine fundiertere Nutzung bestehender Abonnements ermöglicht werden.

Für die praktische Umsetzung wurden mehrere Teilziele verfolgt. Einerseits sollte eine benutzerfreundliche Oberfläche entstehen, die eine klare und verständliche Verwaltung persönlicher Abonnements ermöglicht. Andererseits sollte eine Such- und Detailfunktion für Filme integriert werden, damit Filmdaten, Anbieterinformationen und persönliche Nutzungssituation gemeinsam betrachtet werden können. Darüber hinaus war ein sicheres Authentifizierungs- und Zugriffskonzept erforderlich, um geschützte Bereiche wie die persönliche Abo-Verwaltung datenschutzkonform bereitzustellen.

Die Diplomarbeit verfolgt damit nicht nur das Ziel, eine funktionierende technische Lösung zu entwickeln, sondern auch zu untersuchen, welche fachlichen Anforderungen für eine solche Anwendung relevant sind. Dazu gehören insbesondere Nutzerbedürfnisse, Plattformabdeckung, Benutzerfreundlichkeit, Sicherheitsaspekte und die sinnvolle Integration bestehender externer Dienste wie TMDB und Keycloak.

\section{Nutzen der Anwendung für Endanwender}
Für Endanwender liegt der wesentliche Nutzen von \glqq Mediahub\grqq{} in der zentralen Verwaltung aller Abonnements innerhalb einer einzigen Anwendung. Statt verschiedene Streaming-Plattformen und persönliche Zahlungsinformationen getrennt im Blick behalten zu müssen, erhalten Nutzende eine gemeinsame Übersicht über ihre aktiven Dienste und die damit verbundenen Informationen. Dadurch kann schneller erkannt werden, welche Abonnements tatsächlich genutzt werden und an welchen Stellen unnötige oder doppelte Kosten entstehen.

Ein zusätzlicher Nutzen entsteht durch die Verbindung von persönlicher Abo-Verwaltung mit der Filmsuche. Nutzende können nach Filmen suchen, Detailinformationen aufrufen und gleichzeitig sehen, auf welchen Plattformen ein Film verfügbar ist. Besonders hilfreich ist dabei die Gegenüberstellung von allgemeinen Anbieterinformationen und den bereits gespeicherten eigenen Abonnements. Diese Funktion unterstützt Nutzende dabei, vorhandene Dienste bewusster zu nutzen und Entscheidungen über neue oder bestehende Abonnements besser einzuordnen.

Auch aus Sicht von Benutzerfreundlichkeit und Sicherheit bietet die Anwendung einen konkreten Mehrwert. Geschützte Bereiche werden über Keycloak abgesichert, sodass persönliche Abo-Daten nicht öffentlich zugänglich sind. Gleichzeitig ermöglicht die strukturierte Oberfläche eine verständliche Navigation zwischen öffentlicher Filmsuche und persönlichem Verwaltungsbereich. Die Anwendung richtet sich damit an Nutzende, die ihre Filmnutzung nicht nur konsumieren, sondern auch organisatorisch effizienter verwalten möchten.

\section{Aufbau der Diplomarbeit}
Die vorliegende Diplomarbeit ist so aufgebaut, dass zunächst die fachliche Problemstellung und anschließend die technische Umsetzung der Anwendung schrittweise erläutert werden. Nach der Einleitung werden in Kapitel 2 die Ausgangssituation und die zentralen Herausforderungen bei Streaming-Abonnements beschrieben. Kapitel 3 vertieft den fachlichen Hintergrund und behandelt unter anderem den Nutzen zentraler Abo-Verwaltung, die Bedeutung von Film-Verfügbarkeiten und datenschutzrelevante Anforderungen.

Kapitel 4 stellt die technologischen Grundlagen der Arbeit vor. Dabei werden die eingesetzten Technologien wie Angular, Keycloak, Tailwind CSS, Quarkus und GitHub im Hinblick auf ihre Rolle innerhalb des Projekts erläutert. In Kapitel 5 folgen Architektur und Systemdesign der Anwendung, einschließlich Gesamtarchitektur, API-Integration, Sicherheitskonzept, Datenbankmodell und Backend-Struktur.

Kapitel 6 beschreibt die konkrete Umsetzung der Anwendung mit Fokus auf die Frontend-Funktionen wie Abo-Verwaltung, Filmsuche, Detailansicht und geschützte Bereiche. Anschließend werden in den weiteren Kapiteln Projektverlauf, Ergebnisse, Ausblick und Zusammenfassung dargestellt. Den Abschluss bilden Literaturverzeichnis, Glossar und das Abbildungsverzeichnis.
